% GENERATED FILE. Edit canonical lessons, then run npm run generate:books.
% canonical-source-hash: fnv1a64:39e9590c0db8ee19
% canonical-lessons: MR-C263-gayak, MR-C263-karyakram, MR-C263-abhineta, MR-C263-abhinetri, MR-C263-kavita

\chapter{Singer, Programme, Actor, Actress, Poem}
\label{ch:mr-singer-programme-actor-actress-poem}
\hlchaptermodality{263}

\begin{chapteropening}
\textbf{By the end of this chapter:} \emph{I can say a singer, a programme, an actor, an actress and a poem.}

Complete the last lesson of chapter 263: I can say a singer, a programme, an actor, an actress and a poem.
\end{chapteropening}

\section[\texorpdfstring{\mr{गायक} / \mr{गायिका} (gāyak / gāyikā)}{गायक / गायिका (gāyak / gāyikā)}]{\mr{गायक} / \mr{गायिका} (gāyak / gāyikā) — a singer}

\label{lesson:MR-C263-gayak}



\begin{quote}
Before the new one: say the Marathi for the government, then the Marathi for a law.
\end{quote}

\subsection*{You'll want to know: \mr{गायक} / \mr{गायिका}}
\textbf{\mr{गायक} / \mr{गायिका}} — \emph{gāyak / gāyikā} — \textquotedblleft{}a singer\textquotedblright{}.

\begin{grammarlens}[title={What you've built}]
One more word for this chapter.
\end{grammarlens}

\subsection*{Your turn}
\begin{itemize}
\raggedright
  \item \emph{Say it:} \emph{gāyak / gāyikā}
  \item \emph{Say it:} \emph{gāyak / gāyikā}, once more
  \item \emph{Say it:} say \emph{gāyak / gāyikā}
  \item \emph{Recall:} say the Marathi for the government, then the Marathi for a law, then say \emph{gāyak / gāyikā} again
\end{itemize}

\begin{tcolorbox}[breakable,colback=teal!4,colframe=teal!35!black,title={Before you move on}]
What does \mr{गायक} / \mr{गायिका} mean? (\textquotedblleft{}A singer\textquotedblright{}.) Say it once more.
\end{tcolorbox}
\section[\texorpdfstring{\mr{कार्यक्रम} (kāryakram)}{कार्यक्रम (kāryakram)}]{\mr{कार्यक्रम} (kāryakram) — a programme, a show}

\label{lesson:MR-C263-karyakram}



\begin{quote}
Before the new one: say the Marathi for a law, then the Marathi for a singer.
\end{quote}

\subsection*{You'll want to know: \mr{कार्यक्रम}}
\textbf{\mr{कार्यक्रम}} — \emph{kāryakram} — \textquotedblleft{}a programme, a show\textquotedblright{}.

\begin{grammarlens}[title={What you've built}]
One more word for this chapter.
\end{grammarlens}

\subsection*{Your turn}
\begin{itemize}
\raggedright
  \item \emph{Say it:} \emph{kāryakram}
  \item \emph{Say it:} \emph{kāryakram}, once more
  \item \emph{Say it:} say \emph{kāryakram}
  \item \emph{Recall:} say the Marathi for a law, then the Marathi for a singer, then say \emph{kāryakram} again
\end{itemize}

\begin{tcolorbox}[breakable,colback=teal!4,colframe=teal!35!black,title={Before you move on}]
What does \mr{कार्यक्रम} mean? (\textquotedblleft{}A programme, a show\textquotedblright{}.) Say it once more.
\end{tcolorbox}
\section[\texorpdfstring{\mr{अभिनेता} (abhinetā)}{अभिनेता (abhinetā)}]{\mr{अभिनेता} (abhinetā) — an actor}

\label{lesson:MR-C263-abhineta}



\begin{quote}
Before the new one: say the Marathi for a singer, then the Marathi for a programme.
\end{quote}

\subsection*{You'll want to know: \mr{अभिनेता}}
\textbf{\mr{अभिनेता}} — \emph{abhinetā} — \textquotedblleft{}an actor\textquotedblright{}.

\begin{grammarlens}[title={What you've built}]
One more word for this chapter.
\end{grammarlens}

\subsection*{Your turn}
\begin{itemize}
\raggedright
  \item \emph{Say it:} \emph{abhinetā}
  \item \emph{Say it:} \emph{abhinetā}, once more
  \item \emph{Say it:} say \emph{abhinetā}
  \item \emph{Recall:} say the Marathi for a singer, then the Marathi for a programme, then say \emph{abhinetā} again
\end{itemize}

\begin{tcolorbox}[breakable,colback=teal!4,colframe=teal!35!black,title={Before you move on}]
What does \mr{अभिनेता} mean? (\textquotedblleft{}An actor\textquotedblright{}.) Say it once more.
\end{tcolorbox}
\section[\texorpdfstring{\mr{अभिनेत्री} (abhinetrī)}{अभिनेत्री (abhinetrī)}]{\mr{अभिनेत्री} (abhinetrī) — an actress}

\label{lesson:MR-C263-abhinetri}



\begin{quote}
Before the new one: say the Marathi for a programme, then the Marathi for an actor.
\end{quote}

\subsection*{You'll want to know: \mr{अभिनेत्री}}
\textbf{\mr{अभिनेत्री}} — \emph{abhinetrī} — \textquotedblleft{}an actress\textquotedblright{}.

\begin{grammarlens}[title={What you've built}]
One more word for this chapter.
\end{grammarlens}

\subsection*{Your turn}
\begin{itemize}
\raggedright
  \item \emph{Say it:} \emph{abhinetrī}
  \item \emph{Say it:} \emph{abhinetrī}, once more
  \item \emph{Say it:} say \emph{abhinetrī}
  \item \emph{Recall:} say the Marathi for a programme, then the Marathi for an actor, then say \emph{abhinetrī} again
\end{itemize}

\begin{tcolorbox}[breakable,colback=teal!4,colframe=teal!35!black,title={Before you move on}]
What does \mr{अभिनेत्री} mean? (\textquotedblleft{}An actress\textquotedblright{}.) Say it once more.
\end{tcolorbox}
\section[\texorpdfstring{\mr{कविता} (kavitā)}{कविता (kavitā)}]{\mr{कविता} (kavitā) — a poem}

\label{lesson:MR-C263-kavita}



\begin{quote}
Before the new one: say the Marathi for an actor, then the Marathi for an actress.
\end{quote}

\subsection*{You'll want to know: \mr{कविता}}
\textbf{\mr{कविता}} — \emph{kavitā} — \textquotedblleft{}a poem\textquotedblright{}.

\begin{grammarlens}[title={What you've built}]
One more word for this chapter.
\end{grammarlens}

\subsection*{Your turn}
\begin{itemize}
\raggedright
  \item \emph{Say it:} \emph{kavitā}
  \item \emph{Say it:} \emph{kavitā}, once more
  \item \emph{Say it:} say \emph{kavitā}
  \item \emph{Recall:} say the Marathi for an actor, then the Marathi for an actress, then say \emph{kavitā} again
\end{itemize}

\begin{tcolorbox}[breakable,colback=teal!4,colframe=teal!35!black,title={Before you move on}]
What does \mr{कविता} mean? (\textquotedblleft{}A poem\textquotedblright{}.) Say it once more.
\end{tcolorbox}
